\documentclass[11pt,letterpaper]{article}
\usepackage[T1]{fontenc}
\usepackage[utf8]{inputenc}
\usepackage{lmodern}
\usepackage{amsmath,amssymb,amsthm,mathtools,mathrsfs}
\usepackage{geometry}
\usepackage{microtype}
\usepackage{booktabs,longtable,array}
\usepackage{enumitem}
\usepackage{xcolor}
\usepackage{hyperref}
\usepackage{fancyhdr}
\geometry{margin=1in}
\hypersetup{colorlinks=true,linkcolor=blue!55!black,urlcolor=blue!55!black,citecolor=blue!55!black,pdftitle={Elementary Mathematical Foundations of Wave Confinement Theory},pdfauthor={Richard J. Reyes}}
\newtheorem{theorem}{Theorem}[section]
\newtheorem{lemma}[theorem]{Lemma}
\newtheorem{corollary}[theorem]{Corollary}
\theoremstyle{definition}
\newtheorem{definition}[theorem]{Definition}
\theoremstyle{remark}
\newtheorem{remark}[theorem]{Remark}
\newcommand{\R}{\mathbb R}
\newcommand{\C}{\mathbb C}
\newcommand{\Z}{\mathbb Z}
\newcommand{\N}{\mathbb N}
\newcommand{\T}{\mathbb T}
\newcommand{\Om}{\Omega}
\newcommand{\norm}[1]{\left\lVert#1\right\rVert}
\newcommand{\abs}[1]{\left\lvert#1\right\rvert}
\newcommand{\inner}[2]{\left\langle#1,#2\right\rangle}
\newcommand{\id}[1]{\textnormal{\textbf{[#1]}}}
\newcommand{\scope}{\par\smallskip\noindent\textbf{Limitation.}\ }
\newcommand{\note}{\par\smallskip\noindent\textbf{Interpretation.}\ }
\setlist{itemsep=3pt,topsep=4pt}
\pagestyle{fancy}
\fancyhf{}
\fancyhead[L]{\small WCT elementary mathematical foundations}
\fancyhead[R]{\small R. J. Reyes}
\fancyfoot[C]{\thepage}
\renewcommand{\headrulewidth}{0.4pt}
\setlength{\headheight}{14pt}
\setlength{\parskip}{0.33em}
\setlength{\parindent}{0pt}
\title{\textbf{Elementary Mathematical Foundations of Wave Confinement Theory}\\[.45em]
\large Proof Compendium for the Corrected WCT Equation Registry}
\author{Richard J. Reyes}
\date{October 26, 2026}
\begin{document}
\maketitle
\begin{abstract}
This note collects complete proofs of standard or elementary statements associated with
25 identifiers in the corrected Wave Confinement Theory (WCT) equation registry:
Sobolev embeddings and counterexamples; finite entropy and counting bounds;
operator and logarithmic-flow identities; calculus of variations, dissipation,
and phase-coherence estimates; and flat-torus spectral computations.
Some counting and physical interpretations require independent hypotheses.
Canonical statements are distinguished from auxiliary identities and the historical
uncorrected formulas. The note proves neither a generic physical confinement theorem,
any standard complexity-class separation, nor absence of higher-derivative ghosts.
\end{abstract}
\tableofcontents
\newpage

\section{Conventions}
\subsection{Fourier transform on $\R^n$}
We use the unitary Fourier transform, as in the Fourier-shell manuscript
\cite{reyesMath}:
\begin{equation}\label{eq:fourier}
\widehat f(\xi)=(2\pi)^{-n/2}\int_{\R^n}e^{-ix\cdot\xi}f(x)\,dx,
\qquad
f(x)=(2\pi)^{-n/2}\int_{\R^n}e^{ix\cdot\xi}\widehat f(\xi)\,d\xi.
\end{equation}
Initially these hold for Schwartz functions; Fourier inversion is understood
in the appropriate $L^2$, distributional, or absolutely convergent sense.
Plancherel's theorem and the multiplier convention give
\begin{equation}
\norm f_{L^2(\R^n)}=\norm{\widehat f}_{L^2(\R^n)},\qquad
\widehat{\partial_j f}=i\xi_j\widehat f,\qquad
\widehat{\Delta f}=-\abs\xi^2\widehat f.
\end{equation}
\subsection{Fourier coefficients on the flat torus}
On $\T^n=(\R/2\pi\Z)^n$ we use the \emph{coefficient} convention of the
E62 manuscript \cite{reyesMath}:
\begin{equation}\label{eq:torusfourier}
\widehat f_k=(2\pi)^{-n}\int_{\T^n}f(x)e^{-ik\cdot x}\,dx,
\qquad f(x)=\sum_{k\in\Z^n}\widehat f_ke^{ik\cdot x},
\end{equation}
so Parseval's identity is
\begin{equation}
\norm f_{L^2(\T^n)}^2=(2\pi)^n\sum_{k\in\Z^n}\abs{\widehat f_k}^2.
\end{equation}
On $\prod_{j=1}^n(\R/L_j\Z)$, use the phases
$\exp(2\pi i\sum_jm_jx_j/L_j)$ and normalized coefficients
$\bigl(\prod_jL_j\bigr)^{-1}\int f\,\overline{e_m}$.
\subsection{Sobolev spaces}
For $s\in\R$, set
\begin{equation}\label{eq:Hs}
H^s(\R^n)=\left\{f\in\mathscr S'(\R^n):
\norm f_{H^s}^2:=\int_{\R^n}(1+\abs\xi^2)^s
\abs{\widehat f(\xi)}^2\,d\xi<\infty\right\}.
\end{equation}
For integer $s=2$ this agrees, up to equivalence of norms, with the
weak-derivative definition involving derivatives of order at most two.
$C_c^\infty$ denotes smooth compactly supported functions.

\subsection{The canonical corrected curvature operator}
In accordance with M2 and E17 of the corrected registry \cite{registry},
for $\varepsilon>0$, $\alpha\in\R$, $z\in\C$ put
\begin{equation}\label{eq:regularizer}
D_\varepsilon(z)=\abs z^2+\varepsilon^2e^{-2\alpha\abs z^2},\qquad
R_\varepsilon(z)=\frac{\overline z}{D_\varepsilon(z)},\qquad
\Theta_\varepsilon[\psi]=-(\Delta\psi)R_\varepsilon(\psi).
\end{equation}
The denominator is strictly positive for all complex $z$. For fixed $z\ne0$,
$R_\varepsilon(z)\to\overline z/\abs z^2=1/z$ as $\varepsilon\downarrow0$.
The earlier additive-denominator expression is a different, potentially
singular object and will not be substituted for \eqref{eq:regularizer}.

\subsection{Variation and regularity conventions}
All displayed integrations by parts are justified by compactly supported
variations, periodic boundary conditions, or explicitly imposed vanishing
boundary terms. For real-valued energies of complex fields, variations are
performed in the real and imaginary parts. In a Hilbert-space gradient
flow, the gradient always refers to the designated Hilbert inner product;
$L^2$ and $H^2$ gradients are generally distinct.

\section{Sobolev embeddings}
\begin{theorem}[\id{M4/E24} Sharp $H^2$ threshold]\label{thm:sobolev}
For integer $n\ge1$,
\begin{equation}\label{eq:threshold}
H^2(\R^n)\hookrightarrow L^\infty(\R^n)
\quad\Longleftrightarrow\quad n\le3.
\end{equation}
The continuous embedding means $\norm f_\infty\le C_n\norm f_{H^2}$
for every $f\in H^2(\R^n)$.
\end{theorem}
\begin{proof}
Assume first $n<4$. Fourier inversion and Cauchy--Schwarz give,
for Schwartz $f$,
\begin{align*}
\abs{f(x)}
&\le(2\pi)^{-n/2}\int_{\R^n}\abs{\widehat f(\xi)}\,d\xi\\
&=(2\pi)^{-n/2}\int_{\R^n}(1+\abs\xi^2)^{-1}
(1+\abs\xi^2)\abs{\widehat f(\xi)}\,d\xi\\
&\le (2\pi)^{-n/2}\left(\int_{\R^n}(1+\abs\xi^2)^{-2}\,d\xi\right)^{1/2}
\norm f_{H^2}.
\end{align*}
To determine the integrability of the first factor, polar coordinates yield
\begin{equation*}
I_n=\abs{\mathbb S^{n-1}}\int_0^\infty
\frac{r^{n-1}}{(1+r^2)^2}\,dr.
\end{equation*}
Near the origin the integrand is $O(r^{n-1})$, hence integrable for
$n\ge1$. At infinity it is asymptotic to $r^{n-5}$, whose integral
converges exactly when $n-5<-1$, equivalently $n<4$.
Therefore $\norm f_\infty\le(2\pi)^{-n/2}I_n^{1/2}\norm f_{H^2}$.
Extension by density gives the result for every $f\in H^2$, with a
bounded continuous representative. The converse for $n\ge4$ is proved
by Theorem~\ref{thm:failure}.
\end{proof}
\begin{remark}
The theorem concerns the \emph{chosen} $H^2$ functional-analytic route.
It does not show that long-lived confinement in four spatial dimensions is
impossible or that every possible stability mechanism must use this embedding.
On general domains $\Omega$, extension/regularity assumptions on $\Omega$
are needed before transferring the Euclidean statement.
\end{remark}

\begin{theorem}[\id{E67} Explicit failure above three dimensions]\label{thm:failure}
For every integer $n\ge4$ there exists $f\in H^2(\R^n)$ such that
$f\notin L^\infty(\R^n)$.
\end{theorem}
\begin{proof}
For $n=4$, choose a radial $\chi\in C_c^\infty(\R^4)$ equal to one on
$\abs x\le e^{-3}$ and zero on $\abs x\ge e^{-2}$. For $x\ne0$ put
\begin{equation*}
f(x)=\chi(x)\log\log\frac e{\abs x}.
\end{equation*}
Near the origin, writing $r=\abs x$ and $L(r)=\log(e/r)$, we have
\begin{equation*}
f'(r)=-\frac1{rL(r)},\qquad
f''(r)=\frac{L(r)-1}{r^2L(r)^2}.
\end{equation*}
A radial function on $\R^4$ satisfies
$\abs{D^2f}^2=\abs{f''}^2+3\abs{f'/r}^2$.
Consequently, for $r$ sufficiently small,
\begin{equation*}
\abs{D^2f(r)}^2\le\frac C{r^4L(r)^2}
\end{equation*}
and hence
\begin{align*}
\int_{B_{r_0}}\abs{D^2f}^2\,dx
&\le C'\int_0^{r_0}\frac{dr}{r[\log(e/r)]^2}\\
&=C'\int_{\log(e/r_0)}^\infty\frac{dy}{y^2}<\infty.
\end{align*}
The $L^2$ integrals of $f$ and $\nabla f$ converge by weaker
singularity estimates. The cutoff terms are smooth and supported away
from the origin. Integration by parts against test functions on punctured
balls, followed by removal of the inner boundary, shows that these
classical derivatives agree with the distributional derivatives:
the respective boundary terms vanish. Thus $f\in H^2(\R^4)$.
But $f(r)\to+\infty$ as $r\downarrow0$, so $f\notin L^\infty$.

For $n\ge5$, choose
\begin{equation*}
0<\alpha_0<\frac n2-2
\end{equation*}
and let $f(x)=\chi(x)\abs x^{-\alpha_0}$ with $\chi$ compactly
supported and equal to one near the origin. For $j=0,1,2$,
\begin{equation*}
\abs{D^jf(x)}\le C_j\abs x^{-\alpha_0-j}
\end{equation*}
near zero, so
\begin{equation*}
\int_{B_{r_0}}\abs{D^jf}^2\,dx\le
C\int_0^{r_0}r^{n-1-2\alpha_0-2j}\,dr<\infty.
\end{equation*}
The last inequality follows from $n-2\alpha_0-4>0$. As above,
the cutoff singularity has the stated weak derivatives, so $f\in H^2$.
Yet $f(x)\to\infty$ at the origin. This completes both cases.
\end{proof}
\begin{remark}
The existence of these unbounded $H^2$ functions is classical and is
consistent with the sharp critical Sobolev theory; see
Adams--Fournier \cite{adamsfournier}.
\end{remark}

\begin{definition}[\id{E65} Critical Sobolev power]\label{def:critical}
For $n>2$ the $H^1$ critical Lebesgue exponent and associated semilinear
power are
\begin{equation}\label{eq:crit}
2^*=\frac{2n}{n-2},\qquad
p_c(n):=2^*-1=\frac{n+2}{n-2}.
\end{equation}
The last equality follows by subtracting one. This is a standard exponent
\emph{definition}, not a new existence or stability theorem
\cite{adamsfournier,evans}.
\end{definition}
\begin{remark}
Criticality is model dependent: the displayed $p_c$ is the usual
energy-critical power for $H^1$-based semilinear equations, not a
universally applicable $H^2$-critical exponent.
\end{remark}

\section{Spectral entropy and counting}
\begin{theorem}[\id{E30} Sharp entropy inequalities]\label{thm:entropy}
If $P=(P_1,\ldots,P_K)$, with $K\ge1$, is a probability vector,
$P_j\ge0$ and $\sum_jP_j=1$, then for
$H(P)=-\sum_jP_j\ln P_j$, with $0\ln0=0$, one has
\begin{equation}\label{eq:entropy}
0\le H(P)\le\ln K.
\end{equation}
The upper equality holds exactly for $P_j=1/K$ for every $j$; the lower
equality holds exactly for a point mass.
\end{theorem}
\begin{proof}
Since $0\le P_j\le1$, every term $-P_j\ln P_j$ is nonnegative.
If the sum is zero, each positive $P_j$ must equal one, and the
normalization forces a point mass. For the upper bound, let
$U_j=1/K$ and use Gibbs' inequality for relative entropy:
\begin{equation*}
0\le D(P\Vert U)=\sum_{j:P_j>0}P_j\ln\frac{P_j}{U_j}
=-H(P)+\ln K.
\end{equation*}
Thus $H(P)\le\ln K$, with equality if and only if $P=U$.
Alternatively this is the strict concavity of the logarithm.
\end{proof}
\scope A finite entropy inequality does not prove monotone entropy
production under an evolution. That requires a separate dynamical law.

\begin{corollary}[\id{E33} Effective-support inequality]\label{cor:support}
If $K_t\ge1$ is the number of positive entries of a discrete
probability distribution, then
\begin{equation}
\boxed{e^{H(t)}\le K_t}.
\end{equation}
Equality holds exactly for a distribution uniform on its support.
\end{corollary}
\begin{proof}
Apply Theorem~\ref{thm:entropy} to the strictly positive entries:
$H(t)\le\ln K_t$. Exponentiate both sides.
\end{proof}
\scope The inequality is a \emph{lower} bound for actual support size;
it is not an upper bound. In particular, support shrinkage is not
established by the entropy inequality alone.

\begin{theorem}[\id{E29} Cumulative exponential pruning]\label{thm:pruning}
Suppose $M_t\ge0$ and $\Delta_t\ge0$ obey
$M_{t+1}\le e^{-\Delta_t}M_t$. Then for every integer $T\ge1$,
\begin{equation}\label{eq:pruning}
\boxed{M_T\le M_0\exp\left(-\sum_{t=0}^{T-1}\Delta_t\right)}.
\end{equation}
\end{theorem}
\begin{proof}
At $T=1$, \eqref{eq:pruning} is exactly the given recurrence.
Assume it holds at $T$. Then
\begin{align*}
M_{T+1}&\le e^{-\Delta_T}M_T\\
&\le e^{-\Delta_T}M_0
\exp\left(-\sum_{t=0}^{T-1}\Delta_t\right)
=M_0\exp\left(-\sum_{t=0}^{T}\Delta_t\right).
\end{align*}
The result follows by induction.
\end{proof}
\begin{corollary}[Fixed fraction]
If $0\le\delta\le1$ and $M_{t+1}\le(1-\delta)M_t$, then
\begin{equation*}
M_T\le(1-\delta)^TM_0.
\end{equation*}
\end{corollary}
\begin{proof}
The claim follows by the same induction. For $\delta<1$ it is
Theorem~\ref{thm:pruning} with $\Delta_t=-\ln(1-\delta)$.
At $\delta=1$, $M_t=0$ for $t\ge1$.
\end{proof}
\scope The \emph{canonical} E29 hypothesis is the exponential factor
$e^{-\Delta_t}$, not an assumed constant retention fraction. Entropy
reduction does not imply the recurrence without a model-specific bridge.

\begin{theorem}[\id{E28} Corrected alpha-drop identity]\label{thm:alpha}
Let $n\in\N$, $m(n)\in\N\cup\{0\}$, $\rho_t(n)\in(0,1]$ and
$\beta(n)\in\R$, and define
\begin{equation}\label{eq:alphadef}
\alpha(n):=1+\frac1n\sum_{t=1}^{m(n)}\log_2\rho_t(n)+\beta(n).
\end{equation}
Then
\begin{equation}\label{eq:alphacriterion}
\alpha(n)<1
\quad\Longleftrightarrow\quad
\beta(n)<-\frac1n\sum_{t=1}^{m(n)}\log_2\rho_t(n).
\end{equation}
\end{theorem}
\begin{proof}
Subtract one from both sides of $\alpha(n)<1$, substitute
\eqref{eq:alphadef}, and rearrange. Each step is reversible.
\end{proof}
\begin{remark}
The retained-fraction factors $\rho_t\in(0,1]$ replace the historical
ratio $(M_t+1)/M_t>1$, which cannot by itself induce an exponent drop.
The formula does not show that a proposed physical or algorithmic filter
actually attains those fractions.
\end{remark}

\begin{theorem}[\id{E32} Uniform exponent margin]\label{thm:margin}
Suppose there exist $n_0\ge1$ and $\eta>0$ such that for every
$n\ge n_0$,
\begin{equation}\label{eq:margin}
\beta(n)+\frac1n\sum_{t=1}^{m(n)}\log_2\rho_t(n)\le-\eta.
\end{equation}
Then
\begin{equation*}
\boxed{\limsup_{n\to\infty}\alpha(n)\le1-\eta<1}.
\end{equation*}
\end{theorem}
\begin{proof}
By \eqref{eq:alphadef} and \eqref{eq:margin},
$\alpha(n)\le1-\eta$ for all $n\ge n_0$. Taking the limit superior
establishes the bound.
\end{proof}
\begin{remark}[Terminology]
A bound of order $2^{(1-\eta)n}$ with $0<\eta<1$ is still
\emph{exponential} in $n$, though its exponent improves on $2^n$.
The usual asymptotic meaning of \emph{subexponential} is $2^{o(n)}$.
The canonical registry title ``Subexponential exploration condition''
should not be interpreted as proving $2^{o(n)}$ from $\limsup\alpha<1$.
\end{remark}

\begin{theorem}[\id{E41} Conditional configuration count]\label{thm:count}
Let an initial finite binary code set $B_0(n)$ satisfy
$\abs{B_0(n)}\le2^n$. Assume a sequence of finite retained sets $B_t(n)$
satisfies
\begin{equation}\label{eq:retain}
\abs{B_t(n)}\le\rho_t(n)\abs{B_{t-1}(n)},\qquad
\rho_t(n)\in(0,1].
\end{equation}
Suppose a finite code set $A_n$ has cardinality
\begin{equation}\label{eq:overhead}
\abs{A_n}\le2^{n\beta(n)}\abs{B_{m(n)}(n)}
\end{equation}
and the physically defined set $C_{\mathrm{curv}}(n)$ admits an
injection into $A_n$. Then
\begin{equation}\label{eq:count}
\boxed{\abs{C_{\mathrm{curv}}(n)}\le2^{n\alpha(n)}}.
\end{equation}
\end{theorem}
\begin{proof}
Repeated application of \eqref{eq:retain} gives
$\abs{B_m}\le\abs{B_0}\prod_{t=1}^m\rho_t(n)$.
An injection cannot increase cardinality, so
\begin{align*}
\abs{C_{\mathrm{curv}}(n)}
&\le\abs{A_n}
\le2^{n\beta(n)}\abs{B_0(n)}\prod_{t=1}^{m(n)}\rho_t(n)\\
&\le2^{n+n\beta(n)+\sum_{t=1}^{m(n)}\log_2\rho_t(n)}
=2^{n\alpha(n)}.
\end{align*}
\end{proof}
\scope The injections and cardinality retention inequalities are
\emph{hypotheses}, not consequences of E28 alone. If the right side is
less than one, the hypotheses force the discrete retained set to be
empty. No polynomial-time enumeration follows from \eqref{eq:count}.

\begin{corollary}[\id{E72} Effective search-space count]\label{cor:eff}
If the effective search space $S_{\mathrm{eff}}(n)$ admits an injection
into a finite code set $A_n$ satisfying the cardinality hypotheses
of Theorem~\ref{thm:count}, then
\begin{equation*}
\boxed{\abs{S_{\mathrm{eff}}(n)}\le2^{n\alpha(n)}}.
\end{equation*}
\end{corollary}
\begin{proof}
Injectivity yields $\abs{S_{\mathrm{eff}}(n)}\le\abs{A_n}$;
apply the bound on $A_n$ in Theorem~\ref{thm:count}.
\end{proof}
\scope Neither E41 nor E72 proves any separation between standard
complexity classes. In particular, a configuration count alone does not
prove a computationally efficient encoding or pruning algorithm.

\section{Operator identities}
\begin{lemma}[Auxiliary Laplacian-multiplication commutator]\label{lem:lapcomm}
Let $f,v\in C^2(\Omega)$ be scalar-valued. With
$[\Delta,f]v:=\Delta(fv)-f\Delta v$,
\begin{equation}\label{eq:lapcomm}
\boxed{[\Delta,f]v=(\Delta f)v+2\nabla f\cdot\nabla v}.
\end{equation}
\end{lemma}
\begin{proof}
Twice applying the product rule yields
\begin{equation*}
\partial_j^2(fv)=(\partial_j^2f)v
+2(\partial_jf)(\partial_jv)+f\partial_j^2v.
\end{equation*}
Sum over $j$ and subtract $f\Delta v$.
\end{proof}
\begin{remark}
This is a useful auxiliary identity, but it is \emph{not} the canonical
E51 operator commutator, which is proved next.
\end{remark}

\begin{theorem}[\id{E51} Frozen-denominator curvature commutator]\label{thm:curvcomm}
Let $D\in C^1(\Omega;\C)$ have no zero, let
$\psi\in C^3(\Omega;\C)$, and define a \emph{linear frozen-coefficient}
operator $\Theta_D[v]:=-(\Delta v)/D$. Applying it componentwise to a
vector field,
\begin{equation}\label{eq:curvcomm}
\boxed{\nabla(\Theta_D[\psi])-\Theta_D[\nabla\psi]
=\frac{\Delta\psi}{D^2}\nabla D}.
\end{equation}
\end{theorem}
\begin{proof}
By the quotient rule,
\begin{equation*}
\nabla\left(-\frac{\Delta\psi}{D}\right)
=-\frac{\nabla\Delta\psi}{D}
+\frac{\Delta\psi}{D^2}\nabla D.
\end{equation*}
Constant-coefficient partial derivatives commute on $\R^n$,
so $\nabla\Delta\psi=\Delta\nabla\psi$. The first term is
$\Theta_D[\nabla\psi]$; subtract it to obtain \eqref{eq:curvcomm}.
\end{proof}
\begin{remark}[Relation to the corrected curvature operator]
The true corrected map $\Theta_\varepsilon[\psi]
=-(\Delta\psi)R_\varepsilon(\psi)$ is nonlinear because its
multiplier depends on $\psi$ and $\overline\psi$.
If one freezes the coefficient at a chosen field and defines
$T_\psi[v]=-R_\varepsilon(\psi)\Delta v$, then the exact related
identity is
\begin{equation*}
\nabla(T_\psi\psi)-T_\psi(\nabla\psi)
=-(\Delta\psi)\nabla(R_\varepsilon(\psi)).
\end{equation*}
Neither formula permits replacing the frozen coefficient by one
evaluated afresh on the vector field $\nabla\psi$.
\end{remark}

\begin{theorem}[\id{EX} Logarithmic Laplacian identity]\label{thm:loglap}
Let $u\in C^2(\Omega;\R)$ and $\psi=e^u>0$. Then
\begin{equation}\label{eq:loglap}
\boxed{\frac{\Delta\psi}{\psi}=\Delta u+\abs{\nabla u}^2}.
\end{equation}
\end{theorem}
\begin{proof}
The chain rule gives $\nabla\psi=e^u\nabla u$.
Therefore
\begin{equation*}
\Delta\psi=\nabla\cdot(e^u\nabla u)
=e^u\abs{\nabla u}^2+e^u\Delta u.
\end{equation*}
Divide by $e^u>0$.
\end{proof}

\begin{theorem}[\id{EY/EZ} Cole--Hopf/logarithmic reduction]\label{thm:colehopf}
Let $u\in C^{2,1}(\Omega\times(0,T);\R)$ satisfy
\begin{equation}\label{eq:logflow}
\partial_tu=\Delta u+\abs{\nabla u}^2.
\end{equation}
Then $\psi=e^u$ satisfies the heat equation
\begin{equation}\label{eq:heat}
\boxed{\partial_t\psi=\Delta\psi}.
\end{equation}
Conversely, every strictly positive $C^{2,1}$ solution of
\eqref{eq:heat} gives a solution of \eqref{eq:logflow} by $u=\ln\psi$.
\end{theorem}
\begin{proof}
Differentiation and Theorem~\ref{thm:loglap} show
\begin{equation*}
\psi_t=e^uu_t=e^u(\Delta u+\abs{\nabla u}^2)=\Delta(e^u).
\end{equation*}
Conversely, positivity permits $u=\ln\psi$ and hence
\begin{equation*}
u_t=\frac{\psi_t}{\psi}
=\frac{\Delta\psi}{\psi}
=\Delta u+\abs{\nabla u}^2.
\end{equation*}
\end{proof}
\begin{remark}
This is the classical logarithmic version of the Cole--Hopf
transformation; see Evans \cite[\S 4.4]{evans}.
It does not globally linearize an arbitrary WCT evolution; the
real logarithm breaks down at zeros or sign changes of the field.
\end{remark}

\begin{theorem}[\id{E9} Polar phase-flux identity]\label{thm:phaseflux}
Suppose $\psi$ is a differentiable complex field and locally, on
$\{\psi\ne0\}$, write $\psi=\sqrt\rho\,e^{i\vartheta}$ with
$\rho=\abs\psi^2$. Then
\begin{equation}\label{eq:phaseflux}
\boxed{\operatorname{Im}(\overline\psi\nabla\psi)
=\rho\nabla\vartheta}.
\end{equation}
\end{theorem}
\begin{proof}
Differentiating the polar representation yields
\begin{equation*}
\nabla\psi=e^{i\vartheta}\nabla\sqrt\rho
+i\sqrt\rho\,e^{i\vartheta}\nabla\vartheta.
\end{equation*}
Multiplying by $\overline\psi=\sqrt\rho\,e^{-i\vartheta}$ gives
\begin{equation*}
\overline\psi\nabla\psi
=\sqrt\rho\,\nabla\sqrt\rho+i\rho\nabla\vartheta.
\end{equation*}
The first term is real, so taking the imaginary part proves
\eqref{eq:phaseflux}.
\end{proof}
\scope This is a kinematic identity. A continuity law
$\rho_t+\nabla\cdot\mathbf S=0$ needs an additional field equation or
conservation principle; it does not follow from \eqref{eq:phaseflux}.

\section{Variational and flow lemmas}
\begin{theorem}[\id{E20} Quadratic-sector positivity]\label{thm:quad}
For $\kappa,\theta,\gamma\in\R$ set
$Q(a,b)=\kappa a^2+\gamma ab+\theta b^2$. Then
\begin{equation}\label{eq:quad}
\boxed{Q(a,b)\ge0\quad\text{for all }a,b\in\R
\iff\kappa\ge0,\ \theta\ge0,\ \gamma^2\le4\kappa\theta}.
\end{equation}
\end{theorem}
\begin{proof}
Taking $(a,b)=(1,0)$ and $(0,1)$ gives $\kappa,\theta\ge0$.
If $\kappa>0$, complete the square:
\begin{equation*}
Q(a,b)=\kappa\left(a+\frac\gamma{2\kappa}b\right)^2
+\left(\theta-\frac{\gamma^2}{4\kappa}\right)b^2.
\end{equation*}
Nonnegativity for every $(a,b)$ requires $\gamma^2\le4\kappa\theta$.
If $\kappa=0$, $Q(a,b)=\gamma ab+\theta b^2$ can be nonnegative for
all $(a,b)$ only if $\gamma=0$, which is the same condition.
Conversely, the completed-square expression is nonnegative whenever
$\kappa>0$ and the three inequalities hold. If $\kappa=0$, those
inequalities imply $\gamma=0$, so $Q=\theta b^2\ge0$.
\end{proof}
\begin{remark}[Sign convention and ghost obstruction]
The canonical E20 sector is
$Q(S,P)=\kappa S^2+\theta P^2-\gamma SP$.
The criterion is unchanged on sending $\gamma\mapsto-\gamma$.
Positivity of this algebraic quadratic sector does \emph{not} establish
boundedness of the canonical Hamiltonian, reduction of higher-time
configuration degrees of freedom, or absence of an Ostrogradsky ghost.
The nondegenerate fixed-background S004 analysis is a separate
canonical-constraint result \cite{s004}.
\end{remark}

\begin{theorem}[\id{E21} Second-derivative Euler--Lagrange equation]\label{thm:EL}
Let
\begin{equation*}
S[u]=\int_\Omega L(u,\partial_\mu u,\partial_\mu\partial_\nu u)\,dx
\end{equation*}
with enough differentiability to justify first variation and
integration by parts twice. If $u$ is stationary for all compactly
supported smooth variations, then
\begin{equation}\label{eq:EL}
\boxed{\frac{\partial L}{\partial u}
-\partial_\mu\!\left(\frac{\partial L}{\partial(\partial_\mu u)}\right)
+\partial_\mu\partial_\nu\!\left(
\frac{\partial L}{\partial(\partial_\mu\partial_\nu u)}\right)=0}.
\end{equation}
\end{theorem}
\begin{proof}
For $v\in C_c^\infty(\Omega)$ consider $u_\varepsilon=u+\varepsilon v$.
The first variation is
\begin{align*}
\left.\frac d{d\varepsilon}S[u_\varepsilon]\right|_{\varepsilon=0}
=\int_\Omega\bigg[
&\frac{\partial L}{\partial u}v
+\frac{\partial L}{\partial(\partial_\mu u)}\partial_\mu v\\
&+\frac{\partial L}{\partial(\partial_\mu\partial_\nu u)}
\partial_\mu\partial_\nu v\bigg]dx.
\end{align*}
Integrating the first derivative term by parts once and the
second-derivative term by parts twice, with vanishing boundary terms,
produces
\begin{equation*}
\delta S[u](v)=\int_\Omega\bigg[
\frac{\partial L}{\partial u}
-\partial_\mu\!\left(\frac{\partial L}{\partial(\partial_\mu u)}\right)
+\partial_\mu\partial_\nu\!\left(
\frac{\partial L}{\partial(\partial_\mu\partial_\nu u)}\right)
\bigg]v\,dx.
\end{equation*}
This vanishes for arbitrary $v$. The fundamental lemma of the calculus
of variations gives \eqref{eq:EL} (distributionally, and classically
when the indicated derivatives exist).
\end{proof}
\scope An Euler--Lagrange equation does not by itself establish that its
Cauchy problem is well posed or ghost free.

\begin{theorem}[\id{TOP3} Hilbert-space gradient dissipation]\label{thm:gradient}
Let $\mathcal H$ be a real Hilbert space, $E:\mathcal H\to\R$
continuously Fr\'echet differentiable, and $u(t)$ a sufficiently
regular solution of
\begin{equation*}
u_t=-\nabla_{\mathcal H}E(u).
\end{equation*}
Then
\begin{equation}\label{eq:gradient}
\boxed{\frac{d}{dt}E(u(t))
=-\norm{\nabla_{\mathcal H}E(u(t))}_{\mathcal H}^2\le0}.
\end{equation}
For all times $T$ on which the solution and chain rule exist,
\begin{equation}\label{eq:gradient-int}
E(u(T))+\int_0^T\norm{\nabla_{\mathcal H}E(u(t))}_{\mathcal H}^2\,dt=E(u(0)).
\end{equation}
\end{theorem}
\begin{proof}
The chain rule and Riesz representation yield
\begin{align*}
\frac{dE}{dt}&=DE(u)[u_t]
=\inner{\nabla_{\mathcal H}E(u)}{u_t}_{\mathcal H}\\
&=-\inner{\nabla_{\mathcal H}E(u)}{\nabla_{\mathcal H}E(u)}_{\mathcal H}
=-\norm{\nabla_{\mathcal H}E(u)}_{\mathcal H}^2.
\end{align*}
Integrating this identity over $[0,T]$ proves
\eqref{eq:gradient-int}.
\end{proof}
\begin{remark}
Canonical TOP3 applies this principle to an energy of curves $\gamma$.
To do that rigorously, one must specify the curve-space metric,
parametrization, regularity, and treatment of self-avoidance.
Energy descent alone does not prove global existence, unique equilibrium,
or convergence to a desired geometry.
\end{remark}

\begin{theorem}[\id{E50} Phase-coherence integrability]\label{thm:coherence}
Let $\psi\in L^2(\Omega;\C)$ and suppose a measurable weak phase
gradient $\nabla\vartheta$ is defined where $\psi\ne0$ and satisfies
$\abs{\nabla\vartheta}\ge\delta>0$ almost everywhere on that region.
Define the integrand to be zero on $\{\psi=0\}$. Then
\begin{equation}\label{eq:coherence}
\mathcal C[\psi]=\int_\Omega\frac{\abs\psi^2}{\abs{\nabla\vartheta}}\,dx,
\qquad
\boxed{0\le\mathcal C[\psi]\le\delta^{-1}\norm\psi_{L^2}^2<\infty}.
\end{equation}
\end{theorem}
\begin{proof}
The gradient-floor assumption gives almost everywhere
\begin{equation*}
0\le\frac{\abs\psi^2}{\abs{\nabla\vartheta}}
\le\frac{\abs\psi^2}{\delta}.
\end{equation*}
Integrating and using $\psi\in L^2$ proves the bound.
\end{proof}
\begin{remark}
Alternatively the globally regularized denominator
$\sqrt{\abs{\nabla\vartheta}^2+\delta^2}$ yields the same estimate
where its gradient is defined. Neither formulation proves that a WCT
flow generates or preserves the required phase-gradient lower bound.
\end{remark}

\begin{theorem}[\id{CM9} First-order mode-system equivalence]\label{thm:firstorder}
Let $\delta_\gamma,\delta_b\in C^2(I)$ and suppose the prescribed
second-order system is
\begin{equation*}
\ddot\delta_\gamma=-c_s^2k^2\delta_\gamma-k^2\Phi,
\qquad
\ddot\delta_b=-\mathcal R c_s^2k^2\delta_\gamma-k^2\Phi.
\end{equation*}
Set $v_\gamma=\dot\delta_\gamma$ and $v_b=\dot\delta_b$.
Then the system is equivalent, with matching initial positions and
velocities, to
\begin{equation}\label{eq:cm9}
\boxed{\begin{aligned}
\dot\delta_\gamma&=v_\gamma,&
\dot v_\gamma&=-c_s^2k^2\delta_\gamma-k^2\Phi,\\
\dot\delta_b&=v_b,&
\dot v_b&=-\mathcal R c_s^2k^2\delta_\gamma-k^2\Phi.
\end{aligned}}
\end{equation}
\end{theorem}
\begin{proof}
Differentiating the definitions of $v_\gamma,v_b$ and using the
second-order equations yields \eqref{eq:cm9}. Conversely,
differentiating the first and third equations in \eqref{eq:cm9} and
substituting the second and fourth recovers the second-order system.
Matching initial data makes these maps inverses.
\end{proof}
\scope This is an ODE equivalence, not a derivation of cosmological
coefficients, gravitational forcing, or predictive physical dynamics.

\section{Flat-torus eigenmodes (CLE5--CLE8)}
\begin{theorem}[\id{CLE5} Flat product-torus eigenfunctions]\label{thm:torusspectrum}
Let
\begin{equation*}
\T^2_{L_1,L_2}=(\R/L_1\Z)\times(\R/L_2\Z),
\qquad L_1,L_2>0,
\end{equation*}
with its flat product metric. For $m,n\in\Z$,
\begin{equation}\label{eq:torusmode}
\varphi_{m,n}(s,t)=\exp\left(2\pi i\left(\frac{ms}{L_1}
+\frac{nt}{L_2}\right)\right)
\end{equation}
is an eigenfunction of $-\Delta$, with
\begin{equation}\label{eq:eigenvalue}
\boxed{\lambda_{m,n}=\left(\frac{2\pi m}{L_1}\right)^2
+\left(\frac{2\pi n}{L_2}\right)^2}.
\end{equation}
The modes constitute a complete orthogonal Fourier basis of $L^2$.
\end{theorem}
\begin{proof}
Direct differentiation gives
\begin{equation*}
\partial_s^2\varphi_{m,n}=-\left(\frac{2\pi m}{L_1}\right)^2\varphi_{m,n},
\quad
\partial_t^2\varphi_{m,n}=-\left(\frac{2\pi n}{L_2}\right)^2\varphi_{m,n}.
\end{equation*}
Since $\Delta=\partial_s^2+\partial_t^2$ for the flat metric,
adding proves \eqref{eq:eigenvalue}. Integer $m,n$ give periodicity,
and completeness is the product Fourier series theorem.
\end{proof}

\begin{lemma}[\id{CLE6} Separation in angular coordinates]\label{lem:separation}
On the flat angular torus with metric $R^2d\theta^2+r^2d\phi^2$,
let $-\Delta\psi=\sigma_\star^2\psi$ and suppose
$\psi(\theta,\phi)=f(\theta)g(\phi)$.
Then, wherever $fg\ne0$,
\begin{equation}\label{eq:separation}
\boxed{\frac{f''}{f}+\frac{R^2}{r^2}\frac{g''}{g}
=-\sigma_\star^2R^2}.
\end{equation}
\end{lemma}
\begin{proof}
The Laplacian is $R^{-2}\partial_\theta^2+r^{-2}\partial_\phi^2$.
Insert $\psi=fg$ into the eigenvalue equation to obtain
$-R^{-2}f''g-r^{-2}fg''=\sigma_\star^2fg$.
Dividing by $fg$ and multiplying by $-R^2$ yields
\eqref{eq:separation}. The undivided equation remains valid at zeros.
\end{proof}

\begin{theorem}[\id{CLE7} Periodic angular ODE]\label{thm:angular}
For integer $m\ge0$, every $2\pi$-periodic classical solution of
$f''+m^2f=0$ is
\begin{equation}\label{eq:angular}
\boxed{f(\theta)=A\cos(m\theta)+B\sin(m\theta)}.
\end{equation}
If $m=0$, periodicity forces the solution to be constant.
\end{theorem}
\begin{proof}
For $m>0$, the characteristic roots are $\pm im$, yielding exactly
the two-dimensional solution space \eqref{eq:angular}; integer $m$
ensures periodicity. For $m=0$, the general solution is $A+B\theta$,
and periodicity requires $B=0$.
\end{proof}

\begin{theorem}[Lowest nonzero flat-torus eigenvalue]\label{thm:lowest}
If $L_1>L_2>0$, the least positive eigenvalue of $-\Delta$ on
$\T^2_{L_1,L_2}$ is
\begin{equation*}
\boxed{\lambda_1=\left(\frac{2\pi}{L_1}\right)^2},
\end{equation*}
attained precisely for $(m,n)=(\pm1,0)$.
\end{theorem}
\begin{proof}
For every nonzero integer pair $(m,n)$, either $m\ne0$ or $n\ne0$.
In the first case,
$\lambda_{m,n}\ge(2\pi/L_1)^2$, with equality only for
$\abs m=1$ and $n=0$. In the second case with $m=0$,
$\lambda_{0,n}\ge(2\pi/L_2)^2>(2\pi/L_1)^2$.
Thus precisely the indicated two modes realize the least positive
eigenvalue.
\end{proof}

\begin{theorem}[\id{CLE8} Fixed-winding spectral selection]\label{thm:winding}
On the flat angular torus with $R,r>0$,
$\psi(\theta,\phi)=Ae^{i\phi}$ is an eigenfunction with eigenvalue
$r^{-2}$. In the sector of separated Fourier modes
$e^{i(m\theta+\phi)}$ with fixed $\phi$-winding $n=1$, this is the
minimum eigenvalue, achieved exactly when $m=0$.
\end{theorem}
\begin{proof}
The flat product spectrum in angular coordinates is
\begin{equation*}
\lambda_{m,n}=\frac{m^2}{R^2}+\frac{n^2}{r^2}.
\end{equation*}
At fixed $n=1$,
\begin{equation*}
\lambda_{m,1}=\frac{m^2}{R^2}+\frac1{r^2}\ge\frac1{r^2},
\end{equation*}
with equality only for $m=0$. The mode $Ae^{i\phi}$ is
exactly $(m,n)=(0,1)$.
\end{proof}
\begin{remark}[Global versus fixed-winding selection]
If $R>r$, then the global lowest nonzero angular mode is
$e^{\pm i\theta}$, with eigenvalue $R^{-2}<r^{-2}$.
Thus canonical CLE8 does not claim the global lowest nonzero
product-torus eigenvalue: its minimization is within a fixed winding
sector. Normalization, chirality, and physical electron identification
require additional assumptions.
\end{remark}

\begin{theorem}[Quantified embedded-torus operator error]\label{thm:embedded}
Let $R>r>0$, and consider the standard torus embedded in $\R^3$:
\begin{equation*}
X(\theta,\phi)=\bigl((R+r\cos\phi)\cos\theta,
(R+r\cos\phi)\sin\theta,r\sin\phi\bigr).
\end{equation*}
Its Laplace--Beltrami operator and the flat product comparison are
\begin{align}
\Delta_g v&=\frac{\partial_\theta^2v}{(R+r\cos\phi)^2}
+\frac{\partial_\phi^2v}{r^2}
-\frac{\sin\phi}{r(R+r\cos\phi)}\partial_\phi v,
\label{eq:embedded}\
\Delta_{\mathrm{flat}}v&=\frac1{R^2}\partial_\theta^2v
+\frac1{r^2}\partial_\phi^2v.
\end{align}
For $v\in C^2(\T^2)$ the absolute operator error satisfies
\begin{equation}\label{eq:toruserror}
\boxed{\begin{aligned}
\norm{(\Delta_g-\Delta_{\mathrm{flat}})v}_{\infty}
\le{}&\frac{2Rr+r^2}{R^2(R-r)^2}\norm{\partial_\theta^2v}_{\infty}\\
&+\frac1{r(R-r)}\norm{\partial_\phi v}_{\infty}.
\end{aligned}}
\end{equation}
\end{theorem}
\begin{proof}
The induced metric is diagonal:
$g_{\theta\theta}=(R+r\cos\phi)^2$,
$g_{\phi\phi}=r^2$, and
$\sqrt{\det g}=r(R+r\cos\phi)$.
Substitute into
\begin{equation*}
\Delta_gv=\frac1{\sqrt{\det g}}
\partial_i\bigl(\sqrt{\det g}\,g^{ij}\partial_jv\bigr)
\end{equation*}
to obtain \eqref{eq:embedded}. Subtracting $\Delta_{\rm flat}$ gives
\begin{align*}
(\Delta_g-\Delta_{\rm flat})v
={}&\left[\frac1{(R+r\cos\phi)^2}-\frac1{R^2}\right]\partial_\theta^2v\\
&-\frac{\sin\phi}{r(R+r\cos\phi)}\partial_\phi v.
\end{align*}
Since $R+r\cos\phi\ge R-r>0$,
\begin{equation*}
\left|\frac1{(R+r\cos\phi)^2}-\frac1{R^2}\right|
\le\frac{2Rr+r^2}{R^2(R-r)^2},
\quad
\left|\frac{\sin\phi}{r(R+r\cos\phi)}\right|
\le\frac1{r(R-r)}.
\end{equation*}
Multiply these bounds by the indicated derivative norms and take the
supremum to obtain \eqref{eq:toruserror}.
\end{proof}
\begin{remark}
The flat product operator is exact only for the \emph{flat product}
metric. On an embedded torus the metric-dependent first derivative
and coefficients must be retained or controlled. The displayed
error is a function-dependent absolute estimate, not a uniform relative
eigenvalue-error theorem, and ``thin torus'' alone should not be treated
as zero error for all families of high-frequency modes.
\end{remark}

\section{Registry cross-reference}
The table distinguishes theorem type and specific WCT source family.
A manuscript reference identifies relevant context, not a claim that
its full physical theory is proved by the elementary lemma.
\footnotesize
\setlength{\tabcolsep}{4pt}
\renewcommand{\arraystretch}{1.20}
\begin{longtable}{@{}p{1.0cm}p{1.92cm}p{3.05cm}p{8.0cm}@{}}
\toprule
\textbf{ID} & \textbf{Result} & \textbf{Status} & \textbf{Supporting full paper or source}\\
\midrule\endhead
M4 & Thm.~\ref{thm:sobolev} & Standard & \emph{Hard Upper Bound on Spatial Dimensionality in WCT} \cite{dimensionality} \\
E24 & Thm.~\ref{thm:sobolev} & Standard & \emph{Hard Upper Bound on Spatial Dimensionality in WCT} \cite{dimensionality} \\
E67 & Thm.~\ref{thm:failure} & Standard counterexample & Dimensionality manuscript; \emph{Sobolev Spaces} \cite{dimensionality,adamsfournier} \\
E65 & Def.~\ref{def:critical} & Standard definition & Dimensionality manuscript and standard Sobolev theory \cite{dimensionality,adamsfournier} \\
E30 & Thm.~\ref{thm:entropy} & Standard & \emph{Self-Emergent Fourier Cymatics} \cite{cymatics} \\
E33 & Cor.~\ref{cor:support} & Elementary & \emph{Self-Emergent Fourier Cymatics} \cite{cymatics} \\
E29 & Thm.~\ref{thm:pruning} & Elementary & Corrected WCT equation registry, entropy/counting section \cite{registry} \\
E28 & Thm.~\ref{thm:alpha} & Elementary & \emph{P vs NP in Curvature-Bounded Wave Computation} \cite{complexity} \\
E32 & Thm.~\ref{thm:margin} & Conditional limit & \emph{P vs NP in Curvature-Bounded Wave Computation} \cite{complexity} \\
E41 & Thm.~\ref{thm:count} & Conditional counting & \emph{P vs NP in Curvature-Bounded Wave Computation} \cite{complexity} \\
E72 & Cor.~\ref{cor:eff} & Conditional counting & \emph{P vs NP in Curvature-Bounded Wave Computation} \cite{complexity} \\
E51 & Thm.~\ref{thm:curvcomm} & Elementary identity & \emph{The Geometry of Resonance}; corrected registry \cite{geometry,registry} \\
EX & Thm.~\ref{thm:loglap} & Standard & \emph{Logarithmic Curvature Flow} \cite{logflow} \\
EY & Thm.~\ref{thm:colehopf} & Standard & \emph{Logarithmic Curvature Flow} \cite{logflow} \\
EZ & Thm.~\ref{thm:colehopf} & Standard & \emph{Logarithmic Curvature Flow} \cite{logflow} \\
E9 & Thm.~\ref{thm:phaseflux} & Elementary identity & \emph{Phase-Flux Field} \cite{pff} \\
E20 & Thm.~\ref{thm:quad} & Elementary algebra & \emph{The Geometry of Resonance}; S004 analysis \cite{geometry,s004} \\
E21 & Thm.~\ref{thm:EL} & Standard variation & \emph{The Geometry of Resonance} \cite{geometry} \\
TOP3 & Thm.~\ref{thm:gradient} & Standard, conditional flow & Corrected WCT registry, topology module \cite{registry} \\
E50 & Thm.~\ref{thm:coherence} & Conditional integral & Corrected WCT registry, coherence module \cite{registry} \\
CM9 & Thm.~\ref{thm:firstorder} & Elementary ODE & Corrected WCT registry, cosmology module \cite{registry} \\
CLE5 & Thms.~\ref{thm:torusspectrum}, \ref{thm:embedded} & Exact/approximate & Corrected WCT registry, curvature-locked electron module \cite{registry} \\
CLE6 & Lem.~\ref{lem:separation} & Standard & Corrected WCT registry, curvature-locked electron module \cite{registry} \\
CLE7 & Thm.~\ref{thm:angular} & Standard & Corrected WCT registry, curvature-locked electron module \cite{registry} \\
CLE8 & Thms.~\ref{thm:lowest}, \ref{thm:winding} & Conditional selection & Corrected WCT registry, curvature-locked electron module \cite{registry} \\
\bottomrule
\end{longtable}
\normalsize

\section{Logical scope and conclusion}
The fully analytic/elementary implications in this compendium are
\begin{enumerate}[label=(\roman*)]
\item the sharp $H^2\to L^\infty$ dimension threshold and explicit
counterexamples above three dimensions;
\item entropy and exponential-iteration bounds, with precise equality
conditions or recurrences;
\item standard differential, logarithmic, phase and commutator identities;
\item quadratic-form positivity, Euler--Lagrange and exact
negative-gradient-flow dissipation formulas; and
\item the exact flat product-torus spectrum, plus an explicit
embedded-torus operator-error estimate.
\end{enumerate}
The exponent-margin conclusion E32 and the discrete bounds E41/E72
require independent uniform retention/counting hypotheses. The phase
coherence estimate E50 requires a phase-gradient floor. None of those
hypotheses is obtained merely by assigning PASS status to a symbolic
registry check.

Crucially, these proofs do \emph{not} show that $n=4$ confinement is
physically impossible; that a positive algebraic quadratic sector avoids
Ostrogradsky ghosts; that entropy control implies actual dynamical
spectral concentration; that a bounded code count implies a complexity
separation; or that an admissible winding-one toroidal mode is a unique
physical particle. Such statements require separate theorems and data.

\begin{thebibliography}{99}
\bibitem{registry}
R.~J. Reyes, \emph{Wave Confinement Theory---Corrected Full Equation Registry},
revision 2.0, June 27, 2026 (with later scope notices).
\url{https://github.com/rickyjreyes/geometry_of_resonance/blob/main/WCT_FULL_EQUATION_LIST_CORRECTED.md}.

\bibitem{reyesMath}
R.~J. Reyes, \emph{Mathematical Manuscripts and Formal Proof Artifacts},
versioned research repository, including E62 and Fourier-shell manuscripts.
\url{https://github.com/rickyjreyes/math}.

\bibitem{adamsfournier}
R.~A. Adams and J.~J.~F. Fournier, \emph{Sobolev Spaces},
2nd ed., Pure and Applied Mathematics 140, Academic Press, 2003.

\bibitem{evans}
L.~C. Evans, \emph{Partial Differential Equations}, 2nd ed.,
Graduate Studies in Mathematics 19, American Mathematical Society, 2010.

\bibitem{dimensionality}
R.~J. Reyes, \emph{Hard Upper Bound on Spatial Dimensionality in
Wave Confinement Theory}, Zenodo, 2025.
\url{https://doi.org/10.5281/zenodo.17081283}.

\bibitem{cymatics}
R.~J. Reyes, \emph{Self-Emergent Fourier Cymatics: Entropic Eigenmodes
out of Chaos}, Zenodo, 2025.
\url{https://doi.org/10.5281/zenodo.17732648}.

\bibitem{complexity}
R.~J. Reyes, \emph{P vs NP in Curvature-Bounded Wave Computation:
A Model-Relative Separation}, Zenodo, 2025.
\url{https://doi.org/10.5281/zenodo.17743607}.

\bibitem{pff}
R.~J. Reyes, \emph{Phase-Flux Field (PFF): Axiomatic Substrate for
Wave Confinement Theory}, Zenodo, 2025.
\url{https://doi.org/10.5281/zenodo.17578766}.

\bibitem{logflow}
R.~J. Reyes, \emph{Logarithmic Curvature Flow, Filament Localization,
and the Geometric Origin of the Lepton Mass Spectrum}, Zenodo, 2026.
\url{https://doi.org/10.5281/zenodo.18936949}.

\bibitem{geometry}
R.~J. Reyes, \emph{The Geometry of Resonance: Wave Confinement Theory
and the Emergence of Mass, Force, and Spacetime}, Zenodo, 2025.
\url{https://doi.org/10.5281/zenodo.15644222}.

\bibitem{s004}
R.~J. Reyes, \emph{Ghost and Constraint Analysis: Corrected S004 Sector},
foundational technical note, 2026.
\url{https://github.com/rickyjreyes/geometry_of_resonance/blob/main/docs/foundational/GHOST_AND_CONSTRAINT_ANALYSIS.md}.
\end{thebibliography}
\end{document}
